\documentclass[10pt,a4paper]{amsart}
\usepackage{amsmath,amssymb,amsthm,graphicx}
\usepackage{verbatim}
\newenvironment{reference}{\comment}{\endcomment}
%\newenvironment{reference}{}{}
\newenvironment{slogan}{\comment}{\endcomment}
\newenvironment{history}{\comment}{\endcomment}
% For commutative diagrams we use Xy-pic
\usepackage[all]{xy}

% We use 2cell for 2-commutative diagrams.
\xyoption{2cell}
% Theorem environments.
%
\theoremstyle{plain}
\newtheorem{theorem}[subsection]{Theorem}
\newtheorem{proposition}[subsection]{Proposition}
\newtheorem{lemma}[subsection]{Lemma}

\theoremstyle{definition}
\newtheorem{definition}[subsection]{Definition}
\newtheorem{example}[subsection]{Example}
\newtheorem{exercise}[subsection]{Exercise}
\newtheorem{situation}[subsection]{Situation}

\theoremstyle{remark}
\newtheorem{remark}[subsection]{Remark}
\newtheorem{remarks}[subsection]{Remarks}

\numberwithin{equation}{subsection}

% Macros
%
\def\lim{\mathop{\mathrm{lim}}\nolimits}
\def\colim{\mathop{\mathrm{colim}}\nolimits}
\def\Spec{\mathop{\mathrm{Spec}}}
\def\Hom{\mathop{\mathrm{Hom}}\nolimits}
\def\Ext{\mathop{\mathrm{Ext}}\nolimits}
\def\SheafHom{\mathop{\mathcal{H}\!\mathit{om}}\nolimits}
\def\SheafExt{\mathop{\mathcal{E}\!\mathit{xt}}\nolimits}
\def\Sch{\mathit{Sch}}
\def\Mor{\mathop{\mathrm{Mor}}\nolimits}
\def\Ob{\mathop{\mathrm{Ob}}\nolimits}
\def\Sh{\mathop{\mathit{Sh}}\nolimits}
\def\NL{\mathop{N\!L}\nolimits}
\def\CH{\mathop{\mathrm{CH}}\nolimits}
\def\proetale{{pro\text{-}\acute{e}tale}}
\def\etale{{\acute{e}tale}}
\def\QCoh{\mathit{QCoh}}
\def\Ker{\mathop{\mathrm{Ker}}}
\def\Im{\mathop{\mathrm{Im}}}
\def\Coker{\mathop{\mathrm{Coker}}}
\def\Coim{\mathop{\mathrm{Coim}}}

% Boxtimes
%
\DeclareMathSymbol{\boxtimes}{\mathbin}{AMSa}{"02}

%
% Macros for moduli stacks/spaces
%
\def\QCohstack{\mathcal{QC}\!\mathit{oh}}
\def\Cohstack{\mathcal{C}\!\mathit{oh}}
\def\Spacesstack{\mathcal{S}\!\mathit{paces}}
\def\Quotfunctor{\mathrm{Quot}}
\def\Hilbfunctor{\mathrm{Hilb}}
\def\Curvesstack{\mathcal{C}\!\mathit{urves}}
\def\Polarizedstack{\mathcal{P}\!\mathit{olarized}}
\def\Complexesstack{\mathcal{C}\!\mathit{omplexes}}
% \Pic is the operator that assigns to X its picard group, usage \Pic(X)
% \Picardstack_{X/B} denotes the Picard stack of X over B
% \Picardfunctor_{X/B} denotes the Picard functor of X over B
\def\Pic{\mathop{\mathrm{Pic}}\nolimits}
\def\Picardstack{\mathcal{P}\!\mathit{ic}}
\def\Picardfunctor{\mathrm{Pic}}
\def\Deformationcategory{\mathcal{D}\!\mathit{ef}}
\usepackage{cleveref}
\crefname{originallemma}{original lemma}{original lemmas}
\Crefname{originallemma}{Original lemma}{Original lemmas}
\makeatletter
\newlabel{lemma-affine-base-change}{{07U8}{source}{}{sourcepointer-1}{}}
\newlabel{lemma-affine-base-change@cref}{{[originallemma][0][]07U8}{[1][0][]source}{}{}{}}
\makeatother
\crefname{originallemma}{original lemma}{original lemmas}
\Crefname{originallemma}{Original lemma}{Original lemmas}
\makeatletter
\newlabel{lemma-affine-vanishing-higher-direct-images}{{073H}{source}{}{sourcepointer-2}{}}
\newlabel{lemma-affine-vanishing-higher-direct-images@cref}{{[originallemma][0][]073H}{[1][0][]source}{}{}{}}
\makeatother
\crefname{originallemma}{original lemma}{original lemmas}
\Crefname{originallemma}{Original lemma}{Original lemmas}
\makeatletter
\newlabel{lemma-vanishing-compute}{{07UZ}{source}{}{sourcepointer-3}{}}
\newlabel{lemma-vanishing-compute@cref}{{[originallemma][0][]07UZ}{[1][0][]source}{}{}{}}
\makeatother
\crefname{originallemma}{original lemma}{original lemmas}
\Crefname{originallemma}{Original lemma}{Original lemmas}
\makeatletter
\newlabel{sites-cohomology-lemma-Leray}{{0732}{source}{}{sourcepointer-4}{}}
\newlabel{sites-cohomology-lemma-Leray@cref}{{[originallemma][0][]0732}{[1][0][]source}{}{}{}}
\makeatother
\crefname{originalreference}{original reference}{original references}
\Crefname{originalreference}{Original reference}{Original references}
\makeatletter
\newlabel{situation-vanishing}{{07UY}{source}{}{sourcepointer-5}{}}
\newlabel{situation-vanishing@cref}{{[originalreference][0][]07UY}{[1][0][]source}{}{}{}}
\makeatother
\crefname{originallemma}{original lemma}{original lemmas}
\Crefname{originallemma}{Original lemma}{Original lemmas}
\makeatletter
\newlabel{spaces-morphisms-lemma-base-change-affine}{{03WI}{source}{}{sourcepointer-6}{}}
\newlabel{spaces-morphisms-lemma-base-change-affine@cref}{{[originallemma][0][]03WI}{[1][0][]source}{}{}{}}
\makeatother
\title{Vanishing along an affine base change in the vanishing situation, with the global sections of the base changed structure sheaf}
\author{The Stacks Project Authors}
\date{stacks-project:spaces-cohomology.tex:local-snapshot-20261012; GNU FDL 1.2 or later}
\begin{document}
\maketitle
\noindent\textit{Editorial note.} All delivered mathematical text is the author's own, copied from the pinned Stacks Project chapter file spaces-cohomology.tex; no statement or proof step was rewritten, reordered or removed.\par\medskip
\noindent\textit{Editorial note.} Preamble provenance: the pinned chapter file opens with \textbackslash{}input\{preamble\} and carries no preamble text of its own, so the project's single Stacks-wide preamble staged at sources/stacks-spaces-cohomology/preamble.tex is excerpted verbatim as three spans, lines 8-12 (verbatim package and the project's own reference/slogan/history display environments), lines 14-18 (Xy-pic setup) and lines 154-217 (theorem declarations and mathematical macros). The wrapper class and the amsmath/amssymb/amsthm/graphicx base packages are supplied by the assembler. No macro was invented or redefined.\par\medskip
\noindent\textit{Editorial note.} Same-chapter pointers are located at their \textbackslash{}label lines inside the pinned chapter file; cross-chapter pointers are located in the project's own tag list staged in the snapshot (sources/stacks-chow/tags.txt), at the verbatim lines carrying their canonical tags.\par\medskip
\noindent\textit{Allowed standard prerequisites (editorial).}\begin{itemize}
\item Vanishing of the higher direct images of an affine morphism, and the Leray spectral sequence. Same chapter spaces-cohomology.tex, label lemma-affine-vanishing-higher-direct-images (invoked at proof line 3053) and Cohomology on Sites, label sites-cohomology-lemma-Leray (invoked at proof line 3054). They transport the vanishing of \$H\textasciicircum{}1\$ from the base to any space affine over it, which is part (1).
\item Affineness is preserved by base change, and morphisms of affine schemes are affine. Morphisms of Spaces, label spaces-morphisms-lemma-base-change-affine (invoked at proof line 3057). Applied to the structure morphism \$X \textbackslash{}to \textbackslash{}Spec(A)\$ and the affine base change \$\textbackslash{}Spec(A') \textbackslash{}to \textbackslash{}Spec(A)\$, it shows that \$X' \textbackslash{}to X\$ is affine.
\item The pushforward of the structure sheaf along an affine base change. Same chapter spaces-cohomology.tex, label lemma-affine-base-change (invoked at proof line 3062). It computes \$(X' \textbackslash{}to X)\_*\textbackslash{}mathcal\{O\}\_\{X'\} = A' \textbackslash{}otimes\_A \textbackslash{}mathcal\{O\}\_X\$, which turns the global sections computation into a tensor product over \$A\$.
\item Global sections of the pullback of a module in the vanishing situation. Same chapter spaces-cohomology.tex, label lemma-vanishing-compute (invoked at proof line 3067). It gives \$\textbackslash{}Gamma(X, A' \textbackslash{}otimes\_A \textbackslash{}mathcal\{O\}\_X) = A'\$, completing the second assertion.
\end{itemize}\medskip
\ifdefined\hypertarget\hypertarget{sourcepointer-1}{}\fi
\noindent\textit{Source reference.} Original lemma 07U8: sources/stacks-spaces-cohomology/pinned-original-source.tex, source lines 1902--1902. Navigation pointer for \textbackslash{}ref\{lemma-affine-base-change\} used inside the counted span (line 3062); the argument that consumes it is retained in the span.\par\medskip
\ifdefined\hypertarget\hypertarget{sourcepointer-2}{}\fi
\noindent\textit{Source reference.} Original lemma 073H: sources/stacks-spaces-cohomology/pinned-original-source.tex, source lines 1522--1522. Navigation pointer for \textbackslash{}ref\{lemma-affine-vanishing-higher-direct-images\} used inside the counted span (line 3053); the argument that consumes it is retained in the span.\par\medskip
\ifdefined\hypertarget\hypertarget{sourcepointer-3}{}\fi
\noindent\textit{Source reference.} Original lemma 07UZ: sources/stacks-spaces-cohomology/pinned-original-source.tex, source lines 3009--3009. Navigation pointer for \textbackslash{}ref\{lemma-vanishing-compute\} used inside the counted span (line 3068); the argument that consumes it is retained in the span.\par\medskip
\ifdefined\hypertarget\hypertarget{sourcepointer-4}{}\fi
\noindent\textit{Source reference.} Original lemma 0732: sources/stacks-chow/tags.txt, source lines 9046--9046. Navigation pointer for \textbackslash{}ref\{sites-cohomology-lemma-Leray\} used inside the counted span (line 3055); the argument that consumes it is retained in the span.\par\medskip
\ifdefined\hypertarget\hypertarget{sourcepointer-5}{}\fi
\noindent\textit{Source reference.} Original reference 07UY: sources/stacks-spaces-cohomology/pinned-original-source.tex, source lines 2991--2991. Navigation pointer for \textbackslash{}ref\{situation-vanishing\} used inside the counted span (line 3042); the argument that consumes it is retained in the span.\par\medskip
\ifdefined\hypertarget\hypertarget{sourcepointer-6}{}\fi
\noindent\textit{Source reference.} Original lemma 03WI: sources/stacks-chow/tags.txt, source lines 5113--5113. Navigation pointer for \textbackslash{}ref\{spaces-morphisms-lemma-base-change-affine\} used inside the counted span (line 3058); the argument that consumes it is retained in the span.\par\medskip
\begin{lemma}
\label{lemma-vanishing-base-change}
In Situation \ref{situation-vanishing}.
\begin{enumerate}
\item Given an affine morphism $X' \to X$ of algebraic spaces, we have
$H^1(X', \mathcal{F}') = 0$ for every quasi-coherent
$\mathcal{O}_{X'}$-module $\mathcal{F}'$.
\item Given an $A$-algebra $A'$ setting $X' = X \times_{\Spec(A)} \Spec(A')$
the morphism $X' \to X$ is affine and $\Gamma(X', \mathcal{O}_{X'}) = A'$.
\end{enumerate}
\end{lemma}

\begin{proof}
Part (1) follows from Lemma \ref{lemma-affine-vanishing-higher-direct-images}
and the Leray spectral sequence (Cohomology on Sites, Lemma
\ref{sites-cohomology-lemma-Leray}). Let $A \to A'$ be as in (2).
Then $X' \to X$ is affine because affine morphisms are preserved under
base change (Morphisms of Spaces, Lemma
\ref{spaces-morphisms-lemma-base-change-affine}) and the
fact that a morphism of affine schemes is affine. The equality
$\Gamma(X', \mathcal{O}_{X'}) = A'$ follows as
$(X' \to X)_*\mathcal{O}_{X'} = A' \otimes_A \mathcal{O}_X$
by Lemma \ref{lemma-affine-base-change} and thus
$$
\Gamma(X', \mathcal{O}_{X'}) =
\Gamma(X, (X' \to X)_*\mathcal{O}_{X'}) =
\Gamma(X, A' \otimes_A \mathcal{O}_X) = A'
$$
by Lemma \ref{lemma-vanishing-compute}.
\end{proof}

\end{document}
